\begin{center}[7.]\par\medskip\textit{Das Induktionsgesetz.}\par\medskip
{\small (Gefunden 1835. Januar 23. Morgens $7^{\mathrm u}$ vor dem Aufstehen.)}
\end{center}
I.\quad Die stromerzeugende Kraft, die an einem Punkt $P$ durch ein Leiterelement $\gamma$ hervorgebracht wird, dessen Abstand von $P$ $=r$ beträgt, ist während des Zeitelements $\mathrm dt$ die Differenz der beiden Werte von $\frac{\gamma}{r}$ zu den Zeitpunkten $t$ und $t+\mathrm dt$, geteilt durch $\mathrm dt$. Dabei ist $\gamma$ nach Größe und Richtung zu berücksichtigen. Der Ausdruck
\[
-\frac{\mathrm d\dfrac{\gamma}{r}}{\mathrm dt}
\]
stellt dies kurz und verständlich dar.

\SourceLeaf{616}II.\quad Die Stärke eines erzeugten Stroms ist
\[
=\frac{\int p\,\mathrm ds}{\theta}
\]
wobei $p$ die stromerzeugende Kraft in jedem Element $\mathrm ds$ des Leiters und $\theta$ der Gesamtwiderstand ist.

\begin{center}[8.]\end{center}
Seien $s,s'$ zwei geschlossene Leiter, $r$ der Abstand zweier Punkte in $s$ und $s'$, $r=\rho\rho$ und $\theta$ der Gesamtwiderstand in $s'$.

Es gelten die folgenden beiden Grundgesetze.

I.\quad Fließen in $s$ und $s'$ galvanische Ströme mit den Intensitäten $\varepsilon,\varepsilon'$, die als positiv gelten, wenn die Ströme die Leiter in der für deren Elemente $\mathrm ds,\mathrm ds'$ gewählten Richtung durchlaufen, so ist die abstoßende Kraft der Elemente
\[
=+\frac{\varepsilon\varepsilon'}{\rho}\,.\,\frac{\mathrm d\mathrm d\rho}{\mathrm ds\,.\,\mathrm ds'}\,.\,\mathrm ds\,.\,\mathrm ds'
\]
oder, wenn man mit $\mathrm d,\mathrm d'$ die partiellen Differentiationen in Bezug auf die beiden Ströme bezeichnet,
\[
=+\frac{\varepsilon\varepsilon'}{\rho}\,.\,\mathrm d\mathrm d'\rho
\]
Diese Kraft wirkt in Richtung der Geraden $r$.

\vspace{3mm}
II.\quad Entsteht während der sehr kurzen Zeit $\delta t$ der Strom in $s$, so tritt dabei an jedem Punkt eine stromerzeugende Kraft der oben erwähnten Art auf. Ihr Maß für das Element $\mathrm ds'$ ist
\[
=-\frac{\varepsilon\,\mathrm ds\,.\,\mathrm ds'}{\delta t\,.\,r}\cos u
\]
wobei $u$ die Neigung der Richtungen $\mathrm ds,\mathrm ds'$ gegeneinander bezeichnet.

Bezeichnet $z$ die Projektion von $r$ auf die Richtung von $\mathrm ds'$, so gilt $\mathrm ds\,.\,\cos u=\mathrm dz$, und die Formel lautet damit
\[
=-\frac{\varepsilon\,\mathrm dz}{\delta t\,.\,r}\,.\,\mathrm ds'
\]
Die gesamte von $is$ in $\mathrm ds'$ erregte stromerzeugende Kraft ist
\SourceLeaf{617}\[
=-\frac{\varepsilon\,.\,\mathrm ds'}{\delta t}\int\frac{\mathrm dz}{r}
\]

Da $\frac{\mathrm dz}{r}-\frac{z\,\mathrm dr}{rr}$ ein vollständiges Differential und sein Integral über den geschlossenen Leiter $s$ daher $=0$ ist, ist das obige Integral auch
\[
=-\frac{\varepsilon\,\mathrm ds'}{\delta t}\int\frac{z\,\mathrm dr}{rr}
\]
oder wegen $-\frac{z}{r}\,.\,\mathrm ds'=\mathrm d'r$
\[
=+\frac{\varepsilon}{\delta t}\int\frac{\mathrm dr\,.\,\mathrm d'r}{r}=+\frac{4\varepsilon}{\delta t}\int\mathrm d\rho\,.\,\mathrm d'\rho
\]
Die gesamte stromerzeugende Kraft ist $=\frac{4\varepsilon}{\delta t}\int\int'\mathrm d\rho\,.\,\mathrm d'\rho$. Folglich hat der Strom in $s'$ während der Zeit $\delta t$ die Intensität
\[
\frac{4\varepsilon}{\theta\delta t}\,.\,\int\int'\mathrm d\rho\,.\,\mathrm d'\rho
\]
wofür man offensichtlich auch
\[
-\frac{4\varepsilon}{\theta\delta t}\,.\,\int\int'\rho\,\mathrm d\mathrm d'\rho
\]
schreiben kann. Es gilt nämlich $\int\,.\,\int'\mathrm d\rho\,.\,\mathrm d'\rho=\int\rho\,\mathrm d\rho-\int\int'\rho\,\mathrm d\mathrm d'\rho$ und $\int\rho\,\mathrm d\rho=0$, da das Integral über die ganze Stromlinie $s$ erstreckt wird.

